% Lösungen Einheit 6 von 6 – Graph und Term (zusammenbau v0.1)
\einheitenkopf{Einheit 6 von 6 · Graph und Term}

\erg{21}{a) nach unten \quad b) $f(-2) = 2$, $f(-1) = -1$, $f(0) = -2$, $f(1) = -1$, $f(2) = 2$ \quad c) Randwerte $f(-2) = 2$ und $f(2{,}5) \approx -8{,}13$ \quad d) $32$ Einheiten auf $8$ cm: $1$ cm $= 4$ Einheiten \quad e) $B(0) = 200$, $B(1) = 300$, $B(2) = 450$, $B(3) = 675$, $B(4) = 1012{,}5$; Hochachse etwa $1$ cm $= 100$ Fische \quad f) B – Nullstellen $0$ und $2$, für $x \to +\infty$ nach unten; A verläuft für große $x$ nach oben, C hat die Nullstelle $2$ nicht \quad g) y-Achse auf der Symmetrieachse (nur gerade Exponenten); x-Achse auf Höhe des Randes, denn $f(0) = -2$ und $f(4) = 0$ \quad h) $f(3) = 2{,}5$; $f'(3) < 0$, weil der Graph dort fällt – der positive Funktionswert sagt nichts über die Steigung \quad i) mindestens $4$: drei Extrempunkte, also Grad mindestens drei plus eins \quad j) $y = 2$ – alle Hochpunkte haben den Wert $2$ \quad k) Randwerte $f(-4) = f(4) = -32$; y-Achse stark gestaucht, etwa $1$ cm $= 10$ Einheiten, damit die Tiefpunkte aufs Blatt passen}

21b)

\wertetabelle[2,-1,-2,-1,2]{x}{f(x)}{-2,-1,0,1,2}\begin{ksys}[xmin=-3,xmax=3,ymin=-3,ymax=3,klein]\funktion{\x^2-2}{f}\end{ksys}


21c)

\begin{ksys}[xmin=-3,xmax=3,ymin=-9,ymax=3,klein]\funktionab{-\x^3+3*\x}{f}{-2}{2.5}\end{ksys}


21e)

\begin{ksys}[xmin=0,xmax=4,ymin=0,ymax=1100,xstep=1,ystep=100,xlabel=t in Jahren,ylabel=B,klein]\funktion{200*1.5^\x}{B}\end{ksys}


21k)

\begin{ksys}[xmin=-4,xmax=4,ymin=-90,ymax=10,xstep=1,ystep=10,klein]\funktion{\x^4-18*\x^2}{f}\end{ksys}

\erg{22}{a) Fehler: der Grad folgt aus den Extrempunkten, nicht aus den Nullstellen. Richtig: drei Extrempunkte, also Grad mindestens $4$ \quad b) An jedem Extrempunkt ist die Ableitung null; die Ableitung hat also mindestens drei Nullstellen und damit mindestens Grad drei, die Funktion selbst mindestens Grad vier. \quad c) $f(-2) = 1$, $f(0) = -1$, $f(2) = 1$ \quad d) $H(6 | 9)$: um 14 Uhr sind mit $9000$ Besuchern die meisten im Park}
